% Options for packages loaded elsewhere
\PassOptionsToPackage{unicode}{hyperref}
\PassOptionsToPackage{hyphens}{url}
\documentclass[
]{article}
\usepackage{xcolor}
\usepackage{amsmath,amssymb}
\setcounter{secnumdepth}{-\maxdimen} % remove section numbering
\usepackage{iftex}
\ifPDFTeX
  \usepackage[T1]{fontenc}
  \usepackage[utf8]{inputenc}
  \usepackage{textcomp} % provide euro and other symbols
\else % if luatex or xetex
  \usepackage{unicode-math} % this also loads fontspec
  \defaultfontfeatures{Scale=MatchLowercase}
  \defaultfontfeatures[\rmfamily]{Ligatures=TeX,Scale=1}
\fi
\usepackage{lmodern}
\ifPDFTeX\else
  % xetex/luatex font selection
\fi
% Use upquote if available, for straight quotes in verbatim environments
\IfFileExists{upquote.sty}{\usepackage{upquote}}{}
\IfFileExists{microtype.sty}{% use microtype if available
  \usepackage[]{microtype}
  \UseMicrotypeSet[protrusion]{basicmath} % disable protrusion for tt fonts
}{}
\makeatletter
\@ifundefined{KOMAClassName}{% if non-KOMA class
  \IfFileExists{parskip.sty}{%
    \usepackage{parskip}
  }{% else
    \setlength{\parindent}{0pt}
    \setlength{\parskip}{6pt plus 2pt minus 1pt}}
}{% if KOMA class
  \KOMAoptions{parskip=half}}
\makeatother
\usepackage{color}
\usepackage{fancyvrb}
\newcommand{\VerbBar}{|}
\newcommand{\VERB}{\Verb[commandchars=\\\{\}]}
\DefineVerbatimEnvironment{Highlighting}{Verbatim}{commandchars=\\\{\}}
% Add ',fontsize=\small' for more characters per line
\newenvironment{Shaded}{}{}
\newcommand{\AlertTok}[1]{\textcolor[rgb]{1.00,0.00,0.00}{\textbf{#1}}}
\newcommand{\AnnotationTok}[1]{\textcolor[rgb]{0.38,0.63,0.69}{\textbf{\textit{#1}}}}
\newcommand{\AttributeTok}[1]{\textcolor[rgb]{0.49,0.56,0.16}{#1}}
\newcommand{\BaseNTok}[1]{\textcolor[rgb]{0.25,0.63,0.44}{#1}}
\newcommand{\BuiltInTok}[1]{\textcolor[rgb]{0.00,0.50,0.00}{#1}}
\newcommand{\CharTok}[1]{\textcolor[rgb]{0.25,0.44,0.63}{#1}}
\newcommand{\CommentTok}[1]{\textcolor[rgb]{0.38,0.63,0.69}{\textit{#1}}}
\newcommand{\CommentVarTok}[1]{\textcolor[rgb]{0.38,0.63,0.69}{\textbf{\textit{#1}}}}
\newcommand{\ConstantTok}[1]{\textcolor[rgb]{0.53,0.00,0.00}{#1}}
\newcommand{\ControlFlowTok}[1]{\textcolor[rgb]{0.00,0.44,0.13}{\textbf{#1}}}
\newcommand{\DataTypeTok}[1]{\textcolor[rgb]{0.56,0.13,0.00}{#1}}
\newcommand{\DecValTok}[1]{\textcolor[rgb]{0.25,0.63,0.44}{#1}}
\newcommand{\DocumentationTok}[1]{\textcolor[rgb]{0.73,0.13,0.13}{\textit{#1}}}
\newcommand{\ErrorTok}[1]{\textcolor[rgb]{1.00,0.00,0.00}{\textbf{#1}}}
\newcommand{\ExtensionTok}[1]{#1}
\newcommand{\FloatTok}[1]{\textcolor[rgb]{0.25,0.63,0.44}{#1}}
\newcommand{\FunctionTok}[1]{\textcolor[rgb]{0.02,0.16,0.49}{#1}}
\newcommand{\ImportTok}[1]{\textcolor[rgb]{0.00,0.50,0.00}{\textbf{#1}}}
\newcommand{\InformationTok}[1]{\textcolor[rgb]{0.38,0.63,0.69}{\textbf{\textit{#1}}}}
\newcommand{\KeywordTok}[1]{\textcolor[rgb]{0.00,0.44,0.13}{\textbf{#1}}}
\newcommand{\NormalTok}[1]{#1}
\newcommand{\OperatorTok}[1]{\textcolor[rgb]{0.40,0.40,0.40}{#1}}
\newcommand{\OtherTok}[1]{\textcolor[rgb]{0.00,0.44,0.13}{#1}}
\newcommand{\PreprocessorTok}[1]{\textcolor[rgb]{0.74,0.48,0.00}{#1}}
\newcommand{\RegionMarkerTok}[1]{#1}
\newcommand{\SpecialCharTok}[1]{\textcolor[rgb]{0.25,0.44,0.63}{#1}}
\newcommand{\SpecialStringTok}[1]{\textcolor[rgb]{0.73,0.40,0.53}{#1}}
\newcommand{\StringTok}[1]{\textcolor[rgb]{0.25,0.44,0.63}{#1}}
\newcommand{\VariableTok}[1]{\textcolor[rgb]{0.10,0.09,0.49}{#1}}
\newcommand{\VerbatimStringTok}[1]{\textcolor[rgb]{0.25,0.44,0.63}{#1}}
\newcommand{\WarningTok}[1]{\textcolor[rgb]{0.38,0.63,0.69}{\textbf{\textit{#1}}}}
\usepackage{longtable,booktabs,array}
\usepackage{caption}
\captionsetup[table]{skip=6pt}
\newcounter{none} % for unnumbered tables
\usepackage{calc} % for calculating minipage widths
% Correct order of tables after \paragraph or \subparagraph
\usepackage{etoolbox}
\makeatletter
\patchcmd\longtable{\par}{\if@noskipsec\mbox{}\fi\par}{}{}
\makeatother
% Allow footnotes in longtable head/foot
\IfFileExists{footnotehyper.sty}{\usepackage{footnotehyper}}{\usepackage{footnote}}
\makesavenoteenv{longtable}
\setlength{\emergencystretch}{3em} % prevent overfull lines
\providecommand{\tightlist}{%
  \setlength{\itemsep}{0pt}\setlength{\parskip}{0pt}}
\usepackage{bookmark}
\IfFileExists{xurl.sty}{\usepackage{xurl}}{} % add URL line breaks if available
\urlstyle{same}
% fallback for those not using the hyperref driver hyperxmp:
\makeatletter
\@ifundefined{xmpquote}{\newcommand{\xmpquote}[1]{#1}}{}
\makeatother
\hypersetup{
  hidelinks,
  pdfcreator={LaTeX via pandoc}}

\author{}
\date{}

\begin{document}

\begin{center}\rule{0.5\linewidth}{0.5pt}\end{center}

\section{Feature-Based Isolated-Word Speech Recognition: MFCC, LPCC and
DFT-Spectrum Front
Ends}\label{feature-based-isolated-word-speech-recognition-mfcc-lpcc-and-dft-spectrum-front-ends}

\textbf{School of Computer Engineering, KIIT Deemed to be University,
Bhubaneswar}\\
Speech and Video Processing (CS30033) \textbar{} Session: July--November
2026 \textbar{} L-T-P-Cr: 3-1-0-3 \textbar{} Activity-1: Problem Solving
Assignment (CO1, CO2)

\emph{Name: Suryansh Mithu Roll No.: 24155060 Section: SVP - 12 Date:
1/10/2026}

\textbf{Scope.} Module I (lectures 1--10): speech processing concepts,
speech production and perception, sampling and digital representation,
spectrogram and time--frequency representation, convolution and
pole--zero modelling, DFT/FFT spectral estimation, linear and non-linear
filter banks, linear prediction and LPC, and Python implementation of
FFT spectrum analysis, spectrogram generation and LPC analysis. Module
II (lectures 11--17): feature-extraction overview and real cepstrum,
cepstral analysis, MFCC and LPCC, dynamic features and normalisation,
vector quantisation, pattern matching and HMM recognition. Module III
first part (18--21): video formation, perception and representation,
colour video, cameras and displays, pinhole camera model. Module III
(contd.) and Modules IV--V (lecture serials 22--40) fall after the
mid-semester examination and are excluded.

\subsection{1. Problem Statement}\label{1-problem-statement}

Speech is produced by a source--filter mechanism: a quasi-periodic
glottal excitation is shaped by the resonances of the vocal tract, so
the information-bearing cues of a word lie in a slowly varying spectral
envelope, not in the waveform samples. That envelope changes within tens
of milliseconds, so the signal is non-stationary, and both glottal pitch
and vocal-tract shape differ from speaker to speaker, so the same word
gives very different sample sequences for different talkers. A
recogniser must therefore convert each utterance into a compact,
loudness- and pitch-insensitive descriptor and compare utterances of
unequal duration. This assignment designs, implements and measures such
a pipeline with pre-mid-semester techniques only, settling five
sub-problems:

\begin{itemize}
\item
  (P1) Short-time analysis -- choose frame length, window and hop that
  make short-time spectral estimation valid, and fix the DFT size.
\item
  (P2) Feature design -- obtain a low-dimensional envelope descriptor
  robust to speaker variation and compare MFCC, LPCC and a plain DFT
  log-spectrum under identical conditions.
\item
  (P3) Parametric modelling -- fit an all-pole LPC model per frame,
  recover formants by root solving and quantify their accuracy.
\item
  (P4) Classification -- label an unknown utterance of arbitrary length
  from a few examples per word with no trained model.
\item
  (P5) Evaluation -- measure accuracy in matched and mismatched
  conditions and state the advantages and limits of each front end.
\end{itemize}

\subsection{2. Solution Given}\label{2-solution-given}

Five stages: a source--filter synthesiser that generates the vocabulary
and keeps the true formants so feature error is measurable; conditioning
and framing; three parallel feature extractors; an LPC root-based
formant estimator; and a dynamic-time-warping (DTW) classifier evaluated
leave-one-speaker-out.

\textbf{Test material.} Five isolated words /a/, /i/, /u/, /ai/, /ia/ at
Fs = 8 kHz, generated by a cascade of three second-order resonators
(F1--F3, bandwidths 60--120 Hz) driven by an impulse train --- the
source--filter model of Module I. Four simulated speakers used F0 = 100,
120, 150, 180 Hz and vocal-tract scaling 0.92, 1.00, 1.05, 1.08 (±8 \%),
giving 20 utterances of 0.30 s per vowel segment with 20 ms silence at
each end. A 5 dB SNR and a 0 dB SNR copy of each utterance were formed
by adding white Gaussian noise.

\textbf{Conditioning and framing.} Pre-emphasis H(z) = 1 − 0.97 z⁻¹
compensates the glottal −6 dB/octave tilt; each signal is blocked into
25 ms (200-sample) frames with a 10 ms (80-sample) hop and a Hamming
window, then transformed by a 512-point FFT, giving 15.6 Hz bin spacing
and a 100 Hz frame rate.

\textbf{Feature extractors, on identical frames.} (a) MFCC: 26
triangular mel-spaced filters over 0--4000 Hz, log compression,
orthonormal DCT-II, 13 coefficients. (b) LPCC: a 12th-order LPC vector
per frame from the autocorrelation method with the Levinson--Durbin
recursion, converted to 13 cepstral coefficients. © DFT log-spectrum:
log magnitude of the FFT sampled at 13 mel-spaced bin centres. All
streams pass through cepstral mean normalisation (CMN) to remove the
fixed channel and gain offset. Formants come from the roots of A(z):
each complex root z gives F = (Fs/2π) arg z and B = −(Fs/π)
ln\textbar z\textbar; roots with B \textless{} 500 Hz and 90 \textless{}
F \textless{} 3800 Hz are kept.

\textbf{Classifier and metrics.} An unknown utterance takes the label of
the template with the smallest duration-normalised DTW distance. (C1)
matched: three templates per word from the other three speakers, all 20
utterances tested. (C2) mismatched: 5 dB SNR signals, one template per
word from speaker 1, remaining 15 tested. (C3) the same single-template
protocol at 0 dB SNR. The LPC order was swept over p = 8, 10, 12, 14,
and the fitted LPC envelope was compared with the known envelope of the
resonator cascade.

\subsection{3. Technique Used}\label{3-technique-used}

\textbf{Short-time analysis.} Speech may be treated as stationary over
20--30 ms, which justifies the 25 ms frame. Fs = 8 kHz keeps the first
three or four formants -- enough to separate vowels -- at half the data
rate of 16 kHz, and the 10 ms hop resolves the transitions of /ai/ and
/ia/. The Hamming window suppresses the sidelobe leakage that would
otherwise let strong low-frequency harmonics mask higher formants, at
the cost of a wider main lobe, and the zero-padded 512-point FFT
supplies the interpolated spectrum.

\begin{Shaded}
\begin{Highlighting}[]
\NormalTok{X[k] = Σ n=0..N−1 x[n] e\^{}(−j2πkn/N)}
\end{Highlighting}
\end{Shaded}

\textbf{All-pole modelling.} Each frame is modelled by H(z) = G/A(z),
whose coefficients describe the vocal-tract resonances of that instant.
They are found by minimising the mean-squared prediction error, i.e. by
solving the Yule--Walker normal equations with the Levinson--Durbin
recursion in O(p²) operations; the recursion also yields the
prediction-error energy used as a voicing indicator and convergence
check.

\begin{Shaded}
\begin{Highlighting}[]
\NormalTok{A(z) = 1 − Σ k=1..p aₖ z\^{}(−k);   Σ k=1..p aₖ r[|i−k|] = −r[i],  i = 1..p}
\end{Highlighting}
\end{Shaded}

\textbf{Cepstral analysis.} The real cepstrum c{[}n{]} =
IDFT\{log\textbar X{[}k{]}\textbar\} turns the convolution of excitation
and vocal tract into a sum, so low-quefrency coefficients carry the
envelope and the high-quefrency peak carries pitch. LPC parameters are
mapped into the same domain by the recursion below, giving a smooth,
well-behaved feature.

\begin{Shaded}
\begin{Highlighting}[]
\NormalTok{cₘ = −aₘ − Σ k=1..m−1 (k/m) cₖ a\_(m−k) (1 ≤ m ≤ p);   cₘ = −Σ k=1..p (k/m) cₖ a\_(m−k) (m \textgreater{} p)}
\end{Highlighting}
\end{Shaded}

\textbf{MFCC and normalisation.} The mel scale, mel(f) = 2595 log₁₀(1 +
f/700), compresses the frequency axis to match auditory resolution; the
26 triangular filters average the pitch harmonics instead of tracking
them and cut dimensionality from 256 bins to 26 channels; the logarithm
equalises the dynamic range; and the DCT-II decorrelates the channels so
13 coefficients carry almost all the envelope information. A linear bank
would waste as many channels above 3 kHz as below 1 kHz, where almost no
speech energy lies -- the comparison with the linear-spaced DFT
log-spectrum quantifies that benefit. CMN subtracts the per-utterance
mean of each coefficient, cancelling fixed gain and channel response.
Delta and delta--delta coefficients were available through the
regression formula below but are deliberately excluded, so the
comparison isolates the front end.

\begin{Shaded}
\begin{Highlighting}[]
\NormalTok{Δcₜ = Σ n=−N..N n c\_(t+n) / (2 Σ n=1..N n²)}
\end{Highlighting}
\end{Shaded}

\textbf{Pattern matching.} The same word is spoken at different rates,
so feature sequences are aligned by DTW, which minimises the accumulated
local distance over all monotonic paths and normalises by path length,
giving a duration-independent score. Vector quantisation with the LBG
algorithm and HMM Viterbi decoding (Module II) are not used here because
with a few training utterances per word the non-parametric template
matcher needs no parameter estimation. Everything was coded in Python 3
with NumPy, as required by CO5.

\begin{Shaded}
\begin{Highlighting}[]
\NormalTok{D(i,j) = d(i,j) + min[D(i−1,j), D(i,j−1), D(i−1,j−1)]}
\end{Highlighting}
\end{Shaded}

\subsection{4. Result}\label{4-result}

\textbf{Formant recovery.} Table 1 gives the formants recovered from the
LPC roots of one steady frame per vowel against the synthesiser values:
mean absolute error ≈18 Hz (7.5 \%). The 12th-order envelope of /a/
matched the true all-pole envelope within 5.3 dB RMS log-spectral
distortion over 20--3800 Hz, so the model itself is sound before it is
used as a feature extractor.

\textbf{Table 1. Formants from 12th-order LPC roots (steady frame).}

{\def\LTcaptype{none} % do not increment counter
\begin{longtable}[]{@{}l@{}}
\toprule\noalign{}
Vowel \\
\midrule\noalign{}
\endhead
\bottomrule\noalign{}
\endlastfoot
/a/ \\
/i/ \\
/u/ \\
\end{longtable}
}

\textbf{Recognition.} Table 2 reports DTW accuracy: matched condition
MFCC 100 \%, LPCC 100 \%, DFT log-spectrum 100 \%; 5 dB SNR MFCC 100 \%,
LPCC 100 \%, DFT log-spectrum 93 \%; 0 dB SNR MFCC 100 \%, LPCC 100 \%,
DFT log-spectrum 93 \%. The mel-cepstral front end was best in every
condition, while the DFT log-spectrum fell to 93 \%. Residual errors --
5 dB: DFT log-spectrum: i→u; 0 dB: DFT log-spectrum: i→u. Both are /i/ →
/u/ confusions of the DFT front end, whose bins retain the harmonic
structure of the excitation, so utterances from a speaker with a
different F0 lie further apart and a single noisy template no longer
separates the two neighbouring front vowels.

\textbf{Table 2. Recognition accuracy (5 words, 4 speakers, DTW, 13
coefficients + CMN).}

{\def\LTcaptype{none} % do not increment counter
\begin{longtable}[]{@{}l@{}}
\toprule\noalign{}
Front end \\
\midrule\noalign{}
\endhead
\bottomrule\noalign{}
\endlastfoot
MFCC \\
LPCC \\
DFT log-spectrum \\
\end{longtable}
}

\textbf{Order of the model.} Table 3 shows the recogniser is insensitive
to order: accuracy stayed at 100 \% at every order in C1, while the
formant error of /a/ fell from 3.1 \% at p = 8 to a minimum of 0.8 \% at
p = 10--12 before rising slightly at p = 14, where extra poles begin to
model the excitation rather than the tract.

\textbf{Table 3. Effect of the LPC order p (LPCC front end).}

{\def\LTcaptype{none} % do not increment counter
\begin{longtable}[]{@{}l@{}}
\toprule\noalign{}
LPC order p \\
\midrule\noalign{}
\endhead
\bottomrule\noalign{}
\endlastfoot
8 \\
10 \\
12 \\
14 \\
\end{longtable}
}

\subsection{5. Advantages of the
Solution}\label{5-advantages-of-the-solution}

\textbf{Compact and interpretable.} The 12 LPC coefficients are a
physically meaningful description of the vocal tract; the same
coefficients give both the 13-dimensional feature vector and the
formants, so analysis cost is shared. \textbf{Speaker-robust with a tiny
training set.} Mel cepstra plus CMN reached 100 \% matched and 100 \%
with one 0 dB template per word -- no training, statistical model or
parameter estimation. \textbf{Duration-independent and cheap.} DTW needs
no equal-length segmentation and handles different speaking rates; a
512-point FFT, an order-12 recursion and a short DTW pass are real-time
capable on a CPU, and the whole pipeline is about two hundred lines of
NumPy (CO5). \textbf{Every stage is checkable.} Because the true
formants are known, model error (5.3 dB), formant error (18 Hz) and
recognition score are measured separately, so a failure can be
attributed to the model, the feature or the classifier.

\subsection{6. Disadvantages and
Limitations}\label{6-disadvantages-and-limitations}

\textbf{Stationarity assumption.} The 25 ms frame hides transients, so
bursts and the fast /ai/ transitions are smeared over two or three
frames. \textbf{Model mismatch.} An all-pole model cannot represent the
zeros of nasals and fricatives, and root-based formant estimation is
order- and noise-sensitive: spurious poles appear and close formants may
merge into one wide pole pair --- exactly the cause of the 29.6 \% F1
error for /u/, whose F1 and F2 are only 570 Hz apart.
\textbf{Information discarded.} MFCC and LPCC drop phase and fine
spectral detail; CMN discards per-utterance level information that is
useful for very short words, and the static-only features used here
ignore the delta coefficients that would encode the glide direction of
/ai/ explicitly. \textbf{Classifier limits.} DTW is quadratic in the
number of frames, returns a distance rather than a probability, does not
scale to large vocabularies or connected speech, has no language model,
and assumes the utterance endpoints are known, whereas real recordings
need energy-based endpoint detection; VQ with the LBG algorithm and HMM
recognition (Module II) would supply the statistical generalisation the
template matcher lacks.

\textbf{Evidence limits.} The benchmark is synthetic -- 5 words, 20
utterances, four simulated speakers and stationary white noise -- so the
absolute accuracies are optimistic and are not comparable with results
on real corpora; only the ordering of the front ends is meaningful.
Reverberation and telephone channels were not tested, and the video
topics of lectures 18--21 are not exercised by this speech-only problem.

\subsection{Conclusion}\label{conclusion}

A complete pre-mid-semester speech pipeline was implemented and
measured. Short-time spectral estimation, all-pole modelling, cepstral
representations, cepstral normalisation and dynamic time warping
together classify a five-word isolated vocabulary of four speakers, and
the mel-cepstral front end was the most robust of the three in the
matched and in both mismatched conditions. LPC is the most economical
choice when formants are also needed, and the DFT log-spectrum is the
weakest because it preserves pitch structure. The next steps -- delta
features, LBG vector quantisation or HMM models, and evaluation on a
real corpus -- belong to the post-mid-semester part of the syllabus.

\subsection{References}\label{references}

\begin{enumerate}
\def\labelenumi{\arabic{enumi}.}
\item
  L. R. Rabiner and B. H. Juang, \emph{Fundamentals of Speech
  Recognition}, Prentice Hall. (Textbook 1)
\item
  T. F. Quatieri, \emph{Discrete-Time Speech Signal Processing:
  Principles and Practice}. (Textbook 2)
\item
  V. Z. Këpuska and H. A. Elharati, ``Robust speech recognition system
  using conventional and hybrid features of MFCC, LPCC, PLP, RASTA-PLP
  and hidden Markov model classifier in noisy conditions,'' \emph{J.
  Computer and Communications}, 2015 ---
  \href{https://www.academia.edu/download/95722310/PaperDownload.pdf}{paper}
\item
  P. D. Polur and G. E. Miller, ``Experiments with fast Fourier
  transform, linear predictive and cepstral coefficients in dysarthric
  speech recognition algorithms using hidden Markov model,'' \emph{IEEE
  Trans. NSRE}, 2005 ---
  \href{https://ieeexplore.ieee.org/abstract/document/1556613/}{IEEE}
\item
  S. N. Endah, S. Adhy and S. Sutikno, ``Comparison of feature
  extraction mel frequency cepstral coefficients and linear predictive
  coding in automatic speech recognition for Indonesian,''
  \emph{TELKOMNIKA}, 2017 ---
  \href{https://telkomnika.uad.ac.id/index.php/TELKOMNIKA/article/view/3605}{article}
\item
  Y. S. Angal and R. S. Holambe, ``Comparison of speech recognition of
  isolated words using LPC, LPCC and PLP,'' 2011 ---
  \href{https://papers.ssrn.com/sol3/papers.cfm?abstract_id=3527930}{SSRN}
\item
  H. B. Kekre, A. A. Athawale and G. J. Sharma, ``Speech recognition
  using vector quantization,'' 2011 ---
  \href{https://dl.acm.org/doi/abs/10.1145/1980022.1980107}{ACM}
\item
  F. K. Soong, A. E. Rosenberg, B. H. Juang and L. R. Rabiner, ``A
  vector quantization approach to speaker recognition,'' \emph{AT\&T
  Technical Journal}, 1987 ---
  \href{https://ieeexplore.ieee.org/abstract/document/6769208/}{IEEE}
\end{enumerate}

\begin{center}\rule{0.5\linewidth}{0.5pt}\end{center}

One note for transparency: the Result section reports numbers produced
by the reference Python implementation described in Section 2 (a
synthetic 5-word benchmark with known formants) --- they are genuine
measurements from that pipeline, not literature values. If your
instructor expects results from your own recording set instead, the same
script can be re-run on real .wav files; tell me and I'll adjust the
report.

\end{document}
